\documentclass[11pt]{article}

\usepackage[T1]{fontenc}
\usepackage[utf8]{inputenc}
\usepackage{lmodern}
\usepackage{microtype}
\usepackage[a4paper,margin=26mm]{geometry}
\usepackage{amsmath,amssymb,mathtools}
\usepackage{booktabs,tabularx}
\usepackage{graphicx}
\usepackage{float}
\usepackage{enumitem}
\usepackage{xcolor}
\usepackage{hyperref}

\hypersetup{
  colorlinks=true,
  linkcolor=blue!55!black,
  citecolor=blue!55!black,
  urlcolor=blue!55!black,
  pdftitle={Strict separable non-saturation for a seven-state prepare-and-measure witness},
  pdfauthor={Xiaojian Zeng},
  pdfsubject={Quantum information; prepare-and-measure witnesses; separability},
  pdfkeywords={absolutely entangled set, prepare-and-measure, separable bound, product rays, Groebner basis}
}

\newcommand{\Hil}{\mathcal H}
\newcommand{\SEP}{\mathrm{SEP}}
\newcommand{\Tr}{\operatorname{Tr}}
\newcommand{\spec}{\operatorname{spec}}
\newcommand{\ket}[1]{\lvert #1\rangle}
\newcommand{\bra}[1]{\langle #1\rvert}
\newcommand{\proj}[1]{\ket{#1}\!\bra{#1}}

\title{Strict separable non-saturation for a seven-state\\
prepare-and-measure witness}
\author{Xiaojian Zeng\\
\small PBC School of Finance, Tsinghua University, Beijing 100190, China\\
\small \texttt{xiaojianzeng@foxmail.com}}
\date{}

\begin{document}
\maketitle

\begin{abstract}
Absolutely entangled sets can be certified in a dimension-bounded
prepare-and-measure experiment when a correlation witness exceeds its maximum
over separable preparations.  For a recently proposed seven-state two-qubit
ensemble, the target value is $3/4$, whereas alternating optimization found a
separable value $0.74990679$ and no relaxation was known to place the separable
maximum strictly below $3/4$.  We prove strict non-saturation for this fixed
witness.  Equality in the Helstrom, frame-purity, and Cauchy--Schwarz bounds
would require seven pure product rays to reproduce all 21 transition
probabilities of the target ensemble.  The six zero transitions admit 729
assignments to the two local qubit factors.  Exact orthogonality closure leaves
nine classes up to exchange of the factors.  Eight have unit Gr\"obner ideal
over $\mathbb Q(\sqrt7)$; the remaining class forces a local Bloch Gram matrix
with determinant $-27/16$ and is therefore unphysical.  Compactness then gives
$Q_{\rm SEP}<3/4$.  The proof establishes existence of a positive separable
gap, but it neither computes the optimum nor provides a numerical robustness
threshold.
\end{abstract}

\noindent\textbf{Keywords:} absolutely entangled set; prepare-and-measure
witness; separable bound; product rays; exact algebraic computation

\section{Introduction}

Entanglement depends on a tensor-product decomposition of the Hilbert space.
An absolutely entangled set (AES) is a collection of states such that, for
every global change of basis, at least one member of the transformed set is
entangled.  Explicit
minimal examples in two-qubit systems were introduced by Cai \emph{et
al.}~\cite{Cai2021}.  Prepare-and-measure experiments offer an operational,
semi-device-independent route to this set property: the communicated Hilbert
space has a known dimension, while the preparation and measurement devices are
otherwise uncharacterized.  Prepare-and-measure self-testing provides the
general setting for witnesses built from pairwise state
distinguishability~\cite{Navascues2023}.

Patra, Taoutioui, and V\'ertesi recently constructed such witnesses for
extended AES ensembles~\cite{Patra2026}.  Their seven-state example
$\widetilde S_3$ has unrestricted value $3/4$.  A see-saw search found the
feasible separable value $0.74990679$, but the available semidefinite
relaxation did not produce an upper bound below $3/4$.  Thus a separable
strategy saturating the unrestricted value was not excluded.

Here we settle this equality question for the exact seven-state instance.  We
first derive a necessary equality condition that does not assume a Wigner
property: saturation would require seven product rays having all the target
transition probabilities.  We then exclude such rays by an exhaustive exact
calculation followed by one elementary Bloch-Gram obstruction.  The resulting
statement is deliberately qualitative,
\begin{equation}
  Q_{\rm SEP}<\frac34.                                      \label{eq:main}
\end{equation}
It proves that a gap exists, but supplies neither its magnitude nor the value
of $Q_{\rm SEP}$.  Consequently, it does not yet give a finite-precision
experimental threshold.

\section{The fixed witness}
\label{sec:model}

Let $\Hil=\mathbb C^2\otimes\mathbb C^2$, with basis
$(\ket{00},\ket{01},\ket{10},\ket{11})$ and
$\ket{\pm}=(\ket0\pm\ket1)/\sqrt2$.  The target rays are
\begin{align}
 \ket{\psi_1}&=\ket{00},&
 \ket{\psi_2}&=\ket{0+},\nonumber\\
 \ket{\psi_3}&=\ket{+0},&
 \ket{\psi_4}&=\frac{\ket{00}+\ket{11}}{\sqrt2},
                                                        \label{eq:firstfour}\\
 \ket{\psi_5}&=\frac{\ket{01}}{\sqrt6}
             -\sqrt{\frac23}\ket{10}+\frac{\ket{11}}{\sqrt6},\nonumber\\
 \ket{\psi_6}&=\ket-\otimes\ket1,\nonumber\\
 \ket{\psi_7}&=a_1\ket{00}-b_1(\ket{01}+\ket{10}+\ket{11}).
                                                        \label{eq:lastthree}
\end{align}
where
\begin{equation}
 a_1=\sqrt{\frac{7-2\sqrt7}{14}},\qquad
 b_1=\sqrt{\frac{7+2\sqrt7}{42}}.                       \label{eq:ab}
\end{equation}
Writing $P_i=\proj{\psi_i}$, the positive weights are
\begin{equation}
 \alpha_1=\cdots=\alpha_4=\frac1{6+2\sqrt7},\qquad
 \alpha_5=\alpha_6=\frac{\sqrt7-1}{8},\qquad
 \alpha_7=\frac{3\sqrt7-7}{4}.                          \label{eq:weights}
\end{equation}
Direct calculation gives the tight-frame identities
\begin{equation}
  \sum_{i=1}^7\alpha_i=1,
  \qquad \sum_{i=1}^7\alpha_iP_i=\frac{I_4}{4}.          \label{eq:frame}
\end{equation}

For every unordered pair $1\leq j<i\leq7$, Bob has an independent binary
effect $0\preceq E_{ij}\preceq I_4$.  With
\begin{equation}
 c_{ij}=\alpha_i\alpha_j\lVert P_i-P_j\rVert_1,
\end{equation}
the witness is
\begin{equation}
 W(\rho,E)=\sum_{j<i}c_{ij}
       \Tr[(\rho_i-\rho_j)E_{ij}].                       \label{eq:witness}
\end{equation}
The separable benchmark $Q_{\rm SEP}$ is the attained maximum of
\eqref{eq:witness} over seven trace-one states
$\rho_i\in\SEP(2{:}2)$ and the 21 effects.  This is a finite,
dimensionless, single-round optimization.  No common measurement constraint is
imposed between distinct pairs.

Set
\begin{equation}
 t_{ij}=\Tr(P_iP_j)=\lvert\langle\psi_i|\psi_j\rangle\rvert^2.
\end{equation}
For the target preparations, each effect may be chosen as the Helstrom
projector for $P_i-P_j$.  Its contribution is
$2\alpha_i\alpha_j(1-t_{ij})$.  Using \eqref{eq:frame},
\begin{equation}
 2\sum_{j<i}\alpha_i\alpha_j(1-t_{ij})
 =1-\Tr\!\left(\sum_i\alpha_iP_i\right)^2
 =1-\frac14=\frac34.                                    \label{eq:baseline}
\end{equation}
This reproduces the unrestricted target value.

\section{What equality would require}
\label{sec:equality}

\textbf{Proposition 1.}
If a separable strategy attains $W=3/4$, then there are seven pure product
projectors $R_i$ satisfying
\begin{equation}
  \Tr(R_iR_j)=t_{ij}\qquad(1\leq j<i\leq7).              \label{eq:matching}
\end{equation}

\emph{Proof.}
The feasible preparation and effect sets are compact, so a maximizing strategy
exists.  With the effects fixed, the score is linear in each preparation.
Every separable preparation is a convex combination of pure product
projectors; replacing it by a component that maximizes the relevant linear
functional cannot lower the score.  Repeating this for all seven preparations
shows that some optimum has pure product preparations $R_i$.

Put $s_{ij}=\Tr(R_iR_j)$ and $F=\sum_i\alpha_iR_i$.  Optimizing each binary
effect independently gives
\begin{equation}
 W\leq2\sum_{j<i}\alpha_i\alpha_j
       \sqrt{(1-t_{ij})(1-s_{ij})}.                      \label{eq:helstrom}
\end{equation}
The two squared norms appearing in Cauchy--Schwarz are
\begin{align}
 A&=\sum_{j<i}\alpha_i\alpha_j(1-t_{ij})=\frac38,
                                                               \label{eq:A}\\
 B&=\sum_{j<i}\alpha_i\alpha_j(1-s_{ij})
   =\frac{1-\Tr F^2}{2}\leq\frac38.                    \label{eq:B}
\end{align}
Here \eqref{eq:A} follows from \eqref{eq:baseline}; in
\eqref{eq:B}, $F$ is a density operator on a four-dimensional space, so
$\Tr F^2\geq1/4$.  Therefore
\begin{equation}
  W\leq2\sqrt{AB}\leq\frac34.                            \label{eq:cs}
\end{equation}
If $W=3/4$, every inequality in \eqref{eq:helstrom}--\eqref{eq:cs} is an
equality.  Hence $A=B$, and equality in Cauchy--Schwarz makes the two
21-component vectors
\begin{equation}
 \bigl(\sqrt{\alpha_i\alpha_j(1-t_{ij})}\bigr)_{j<i},
 \qquad
 \bigl(\sqrt{\alpha_i\alpha_j(1-s_{ij})}\bigr)_{j<i}
\end{equation}
proportional.  Their norms agree.  All weights are positive and all target
rays are distinct, so the proportionality factor is one and
$s_{ij}=t_{ij}$ for every pair.  This proves \eqref{eq:matching}.
\hfill$\square$

The proposition only uses equality in metric and purity bounds.  In
particular, it does not require the transition probabilities to determine the
target rays up to a unitary or antiunitary, and hence does not invoke a Wigner
property.

\section{Exact product-ray obstruction}
\label{sec:obstruction}

Write a product ray as
$R_i=\proj{a_i}\otimes\proj{b_i}$.  For unit Bloch vectors
$\boldsymbol r_i^A,\boldsymbol r_i^B\in S^2$, define the local fidelities
\begin{equation}
 q_{ij}^X=\lvert\langle x_i|x_j\rangle\rvert^2
 =\frac{1+\boldsymbol r_i^X\!\cdot\boldsymbol r_j^X}{2},
 \qquad X\in\{A,B\}.                                    \label{eq:q}
\end{equation}
Condition \eqref{eq:matching} becomes
\begin{equation}
 q_{ij}^Aq_{ij}^B=t_{ij}.                                \label{eq:factor}
\end{equation}
The complete target matrix is given in Appendix~\ref{app:transitions}.  Its
zero graph is
\begin{equation}
 Z=\{15,16,36,56,57,67\},                                \label{eq:zeros}
\end{equation}
where, for example, $15$ denotes the pair $\{1,5\}$.

We use the following finite exact test.  Each edge in $Z$ is assigned to
factor $A$, factor $B$, or both, according to which local fidelity vanishes.
This gives $3^6=729$ assignments.  For a qubit ray, the orthogonal ray is
unique: in Bloch coordinates an orthogonality edge makes the two vectors
antipodal.  Thus every connected local zero graph has a parity labeling;
equal-parity vertices represent the same ray, while opposite-parity vertices
are orthogonal.  Odd cycles and closures that make a nonzero target edge
orthogonal are impossible.  Exact closure leaves 26 assignments, 18 ordered
closure patterns, and nine patterns up to $A\leftrightarrow B$.

For each remaining pattern, assign a formal dot product to each pair of
distinct parity components.  If vertices $i,j$ have component signs
$\sigma_i,\sigma_j\in\{\pm1\}$, then
\begin{equation}
 q_{ij}^X=\frac{1+\sigma_i\sigma_j x^X_{C(i)C(j)}}{2}.
                                                               \label{eq:componentq}
\end{equation}
Substitution in the 15 nonzero equations \eqref{eq:factor} produces a
polynomial system over $\mathbb Q(\sqrt7)$.  Exact Gr\"obner reduction gives
the unit ideal for eight of the nine exchange classes.  Appendix
\ref{app:enumeration} lists all classes and deletion-minimal contradictory
sets; the accompanying checker prints the factored equations.  Because a
physical Bloch solution would also solve these polynomial equations, a unit
ideal is already a conclusive rejection.  No positivity relaxation is used in
this step.

The sole algebraic survivor, up to exchange of the two factors, has the local
closure data
\begin{align}
 Z_A&=\{15,16,57,67\}, & O_A&=\{17,56\},\nonumber\\
 Z_B&=\{36,56\},       & O_B&=\{35\},                    \label{eq:survivor}
\end{align}
where $O_X$ denotes pairs forced to have local fidelity one.  Reduction by the
surviving transition ideal gives
\begin{equation}
 q^A_{23}=q^A_{34}=1.                                    \label{eq:forcedone}
\end{equation}
Unit Bloch vectors with dot product one coincide.  Hence rays $2,3,4$ in
factor $A$ coincide and $q^A_{24}=1$ as well.  The target equations then force
in factor $B$
\begin{equation}
 q^B_{1k}=\frac34\quad(k=2,3,4),\qquad
 q^B_{k\ell}=\frac14\quad(2\leq k<\ell\leq4).            \label{eq:forcedq}
\end{equation}
The Bloch Gram matrix of these four local rays would therefore be
\begin{equation}
 G=\begin{pmatrix}
 1&1/2&1/2&1/2\\
 1/2&1&-1/2&-1/2\\
 1/2&-1/2&1&-1/2\\
 1/2&-1/2&-1/2&1
 \end{pmatrix}.                                          \label{eq:gram}
\end{equation}
But
\begin{equation}
 \det G=-\frac{27}{16},\qquad
 \spec(G)=\left\{-\frac12,\frac32,\frac32,\frac32\right\}. \label{eq:gramspectrum}
\end{equation}
Every real Gram matrix is positive semidefinite, so this survivor is
impossible.  The factor-exchanged class has the same obstruction.

\textbf{Lemma 2 (computer-assisted product obstruction).}
No seven two-qubit product rays reproduce the 21 target transition
probabilities.

The argument above proves the lemma, with the finite enumeration and exact
ideal calculations supplied as reproducible algebraic certificates.  The only
non-ideal step for the survivor is the elementary unit-vector implication
$\boldsymbol r\cdot\boldsymbol s=1\Rightarrow\boldsymbol r=\boldsymbol s$,
followed by the exact spectrum \eqref{eq:gramspectrum}.

\section{Strictness and physical scope}
\label{sec:strictness}

\textbf{Theorem 3.}
For the states, weights, witness normalization, two-qubit separable
preparations, and independent binary effects of Section~\ref{sec:model},
\begin{equation}
 Q_{\rm SEP}<\frac34.
\end{equation}

\emph{Proof.}
The two-qubit separable-state set and each effect interval
$\{E:0\preceq E\preceq I_4\}$ are compact.  Their finite Cartesian product is
compact, and \eqref{eq:witness} is continuous, so $Q_{\rm SEP}$ is attained.
Equation \eqref{eq:cs} gives $Q_{\rm SEP}\leq3/4$.  If equality held,
Proposition~1 would produce product rays satisfying \eqref{eq:matching}, in
contradiction with Lemma~2.  Therefore the inequality is strict.
\hfill$\square$

Together with the feasible value reported in Ref.~\cite{Patra2026}, the
result changes the known interval to
\begin{equation}
  0.74990679\leq Q_{\rm SEP}<0.75,                        \label{eq:interval}
\end{equation}
where the lower endpoint remains a numerical see-saw value.  The strict upper
endpoint means that the ideal target correlations lie outside the separable
correlation set.  Under the assumed four-dimensional message and the
prepare-and-measure causal model, the fixed witness therefore has a genuine
AES-certification gap for the seven-state extended ensemble.

Several boundaries matter.  The result is not device independent; the
four-dimensional message assumption is essential.  It concerns the extended
seven-state ensemble rather than directly certifying the original four-state
subset with this witness.  It also neither determines $Q_{\rm SEP}$ nor proves
that the see-saw point in \eqref{eq:interval} is globally optimal.  Most
importantly for experiment, compactness supplies no computable modulus: this
work gives no explicit number $U<3/4$ with certified decimal value and hence no
finite noise tolerance.

\section{Methods and reproducibility}
\label{sec:methods}

The target states and weights were independently entered in coordinate and
tensor-product representations before the present calculation.  Both give
unit-normalized rays, positive weights summing to one, the frame identity
\eqref{eq:frame}, legal pairwise Helstrom projectors, and the exact value
\eqref{eq:baseline}.

The product obstruction is implemented in a deterministic Python script using
SymPy~\cite{Meurer2017}.  It first compares all 21 embedded target transitions
against the independently generated machine-readable target input.  It then
enumerates the $3^6$ zero-edge assignments, performs parity closure, constructs
the exact systems \eqref{eq:componentq}, and computes Gr\"obner bases over the
algebraic extension $\mathbb Q(\sqrt7)$.  Rejected systems contain $1$ in the
ideal; the script also deletes equations until it obtains the cores summarized
in Appendix~\ref{app:enumeration}.  The surviving system is reduced exactly to
\eqref{eq:forcedone}--\eqref{eq:forcedq}, and the determinant and characteristic
polynomial of \eqref{eq:gram} are evaluated over the rationals.  No
floating-point comparison or fitted tolerance enters the proof.  The code,
input, exact output, environment specification, and run instructions accompany
the manuscript.

\section{Conclusion}

The seven-state prepare-and-measure witness cannot be saturated by separable
two-qubit preparations.  The central mechanism is rigid: equality would copy
every target transition probability into a product-ray factorization, while
the target's six zero transitions reduce all such factorizations to finitely
many local orthogonality patterns.  Exact algebra eliminates eight patterns;
the last violates positivity of a four-vector Bloch Gram matrix.  This proves
the existential separation $Q_{\rm SEP}<3/4$ without assuming a Wigner
property.

The next quantitative task is to make the obstruction stable.  A bound that
turns a small witness deficit into approximate transition matching, combined
with a robust version of the zero-graph and negative-Gram arguments, could
produce a rational or interval-certified $U<3/4$.  Such a bound, or a tight
moment/SOS dual certificate, is required for an operational experimental
threshold and for deciding whether the reported see-saw value is optimal.

\section*{Statements and declarations}

\noindent\textbf{Data and code availability.}
All target data used in the proof are displayed in this article.  The exact
checker, machine-readable input, output certificate, dependency versions, and
regeneration instructions are supplied as ancillary files with the manuscript.

\medskip
\noindent\textbf{Generative-AI disclosure.}
OpenAI Codex using the GPT-5.6 Sol model assisted with idea development,
symbolic-algebra code, exact calculations, literature searching, verification,
and manuscript drafting.  The AI system is not an author.  The corresponding
author will review the complete scientific and editorial record and remains
responsible for any version that is submitted.

\medskip
\noindent\textbf{Funding and competing interests.}
The final funding and competing-interest statements require confirmation by
the author before submission.

\appendix

\section{Target transition probabilities}
\label{app:transitions}

For completeness, define
\begin{equation}
 u=\frac{7-2\sqrt7}{14},\qquad
 v=\frac{14-5\sqrt7}{42}.
\end{equation}
Direct evaluation of \eqref{eq:firstfour}--\eqref{eq:ab} gives
\begin{equation}
 (t_{ij})_{i,j=1}^7=
 \begin{pmatrix}
 1&1/2&1/2&1/2&0&0&u\\
 1/2&1&1/4&1/4&1/12&1/4&v\\
 1/2&1/4&1&1/4&1/3&0&v\\
 1/2&1/4&1/4&1&1/12&1/4&v\\
 0&1/12&1/3&1/12&1&0&0\\
 0&1/4&0&1/4&0&1&0\\
 u&v&v&v&0&0&1
 \end{pmatrix}.                                          \label{eq:tmatrix}
\end{equation}
Every off-diagonal entry is strictly smaller than one.  The six vanishing
entries in its upper triangle are exactly those in \eqref{eq:zeros}.

\section{Finite enumeration certificate}
\label{app:enumeration}

Table~\ref{tab:classes} records the nine closure classes produced by the exact
parity algorithm.  For each factor, $Z_X$ contains pairs forced orthogonal and
$O_X$ contains pairs forced identical.  The rejected-row entry gives a
deletion-minimal set of target-pair equations whose exact Gr\"obner basis is
$[1]$.  The accompanying output lists the corresponding factored polynomials,
so each rejection can be replayed without a numerical threshold.

\begin{table}[H]
\centering
\scriptsize
\caption{Closure classes up to exchange of local factors.  Pair labels are
one-based and unordered.}
\label{tab:classes}
\begin{tabularx}{\textwidth}{@{}cXXX@{}}
\toprule
Class & $A$: $Z_A/O_A$ & $B$: $Z_B/O_B$ & Exact outcome \\
\midrule
1 & \shortstack[l]{$Z=\{15,16,57,67\}$\\$O=\{17,56\}$}
  & \shortstack[l]{$Z=\{36,56\}$\\$O=\{35\}$}
  & Algebraic survivor; rejected by \eqref{eq:gram}. \\
2 & \shortstack[l]{$Z=\{15,16,57,67\}$\\$O=\{17,56\}$}
  & \shortstack[l]{$Z=\{36,56,67\}$\\$O=\{35,37,57\}$}
  & Unit ideal; core $\{45,46,47\}$. \\
3 & \shortstack[l]{$Z=\{15,36,57\}$\\$O=\{17\}$}
  & \shortstack[l]{$Z=\{16,56,67\}$\\$O=\{15,17,57\}$}
  & Unit ideal; core $\{17\}$. \\
4 & \shortstack[l]{$Z=\{15,56,57\}$\\$O=\{16,17,67\}$}
  & \shortstack[l]{$Z=\{16,36,67\}$\\$O=\{13,17,37\}$}
  & Unit ideal; core $\{17\}$. \\
5 & \shortstack[l]{$Z=\{15,16,57,67\}$\\$O=\{17,56\}$}
  & \shortstack[l]{$Z=\{16,36,56\}$\\$O=\{13,15,35\}$}
  & Unit ideal; core $\{17,35,37\}$. \\
6 & \shortstack[l]{$Z=\{15,57\}$\\$O=\{17\}$}
  & \shortstack[l]{$Z=\{16,36,56,67\}$\\$O=\{13,15,17,35,37,57\}$}
  & Unit ideal; core $\{35,37\}$. \\
7 & \shortstack[l]{$Z=\{15,16,57,67\}$\\$O=\{17,56\}$}
  & \shortstack[l]{$Z=\{16,36,56,67\}$\\$O=\{13,15,17,35,37,57\}$}
  & Unit ideal; core $\{45,46,47\}$. \\
8 & \shortstack[l]{$Z=\{15,56,57\}$\\$O=\{16,17,67\}$}
  & \shortstack[l]{$Z=\{16,36,56,67\}$\\$O=\{13,15,17,35,37,57\}$}
  & Unit ideal; core $\{45,46,47\}$. \\
9 & \shortstack[l]{$Z=\{15,36,57\}$\\$O=\{17\}$}
  & \shortstack[l]{$Z=\{16,36,56,67\}$\\$O=\{13,15,17,35,37,57\}$}
  & Unit ideal; core $\{35,37\}$. \\
\bottomrule
\end{tabularx}
\end{table}

The counts provide additional completeness checks: 26 of the 729 assignments
survive local closure; these produce 18 ordered patterns; every ordered pattern
is one of the nine rows or its factor exchange.  Thus a product-ray solution
cannot lie outside the table.

\bibliographystyle{unsrt}
\bibliography{manuscript/references}

\end{document}
